\section{Wantzel's proofs}

\begin{theorem}
Every two valued function of $n$ variables is of the form $S_1 \pm S_2\sqrt{\triangle}$, where $S_1$ and $S_2$ are integral symmetric functions, and $\triangle$ is the discriminant.
\end{theorem}

\begin{proof}
A two-valued function must belong the the group of order $\half N$. The only group of this order is the alternating group to which the function $\sqrt{\triangle}$ belongs. The theorem therefore 
follows as an immediate consequence of the fundamental theorem Thm (\ref{FunctionsOfSameGroup}). \\

{\bf An Second Proof}\\
Let two values of the function be denoted by $\phi_1$ and $\phi_2$ and let $G_1$ and $G_2$ be the corresponding groups, each of order $\half N$. In the first place, these two groups must
be identical; for it any substitution $S$ of $G_1$ were to change $\phi_2$ to its second value $\phi_1$, then $S^{-1}$ would change $\phi_1$ into $\phi_2$; but this is impossible, since $S^{-1}$ and $S$ each belong to $G_1$. Every substitution therefore of $G_1$ must belong to $G_2$ and vice-versa. \\

To show now that these groups coincided with the alternating group, consider the function $\psi = \phi_1 - \phi_2$. Any substitution which belongs to the common group leaves this unaltered; any other will change $phi_1$ to $\phi_2$ and $\phi_2$ to $\phi_1$, and will therefore change the sign of $\psi$; some transposition, $(x_\alpha x_\beta)$ for example, will have this affect, for no group can include all the transpositions without coinciding with the symmetric group. It is easily inferred that $\psi_1 - \psi_2$ is divisible by $x_\alpha - x_\beta$, and hence by the product of all the differences, since $\psi^2$ is symmetric. \\

The quotient of $\psi$ by $\sqrt{\triangle}$ is symmetric. To prove this, let $(\sqrt{\triangle})^m$ be the highest power of $\sqrt{\triangle}$ which occurs in $\psi$. The quotient of $\psi$ by
$(\sqrt{\triangle}^m$ is symmetric, since, if not, it would be an alternating function, and would again contain $\sqrt{\triangle}$ as a factor, which is contrary to assumption. It follows immediately that $m$ is an odd number and that the quotient of $\psi$ by $\sqrt{\triangle}$ is symmetric. Writing therefore $\psi_1 - \psi_2 = A\sqrt{\triangle}$ and $\phi_1 + \phi_2 = B$, where $A$ and $B$ are both symmetric, we have 
\[ \phi_1 = S_1 + S_2\sqrt{\triangle} \mbox{  and  } \phi_2 = S_1 - S_2\sqrt{\triangle},\] where $S_1$ and $S_2$ are both symmetric functions of the variables $x_1, x_2, \hdots, x_n$. The groups $G_1$ and $G_2$ obviously coincide with the group of $\sqrt{\triangle}$ -- the Alternating Group. 
\end{proof}

\begin{theorem}
Permutations are elements of the Symmetric Group and act on products of functions homomorphically 
because the action naturally distributes over multiplication. 
By definition, a permutation rearranges inputs, which respects function evaluation and standard pointwise multiplication. 
The following operations define the action of the Symmetric Group:
\begin{eqnarray*}
\sigma \cdot (f + g) &=& \sigma\cdot f + \sigma\cdot g;\cr
\sigma\cdot(fg) &=& (\sigma\cdot f)(\sigma \cdot g),\cr
\tau\cdot (\sigma\cdot f) &=& (\tau\sigma)\cdot f
\end{eqnarray*}for $\sigma, \tau \in S_n$ and $f, g\in F[x_1,\hdots, x_n)]$.  $S_n$ is the set of permutations on $n$ elements $[x_1, x_2, \hdots, x_n]$.
\end{theorem}
\begin{proof} 

{\bf 1. Define the Permutation Action A permutation} $\sigma$ acts on a function $f$  by precomposition \cite{Cox}: \[ (\sigma \cdot f)(x) = f(\sigma^{-1}(x)) \] 

{\bf 2. Evaluate the Action on a Product of Functions} We apply the permutation  to the pointwise product of two functions, 
$(fg)(x) = f(x) g(x)$. \[ [\sigma\cdot (fg)](x) = (fg)(\sigma^{-1}(x)).\]

{\bf 3. Unpack the Function Evaluation} Evaluate the product at the permuted input $\sigma^{-1}(x)$:
\[ [\sigma\cdot (fg)](x) = f(\sigma^{-1}(x)) \cdot g(\sigma^{-1}(x)).\]

{\bf 4. Regroup using the Action Definition} By the definition of the action in Step 1, we rewrite the scalar outputs:
\[ [\sigma\cdot (fg)](x) = (\sigma\cdot f)(x) \cdot (\sigma \cdot g)(x).\]

{\bf 5. Conclude for the Pointwise Product} Since this holds for every element $x\in X$, the functions are equivalent:
\[ \sigma\cdot(fg) = (\sigma\cdot f)(\sigma \cdot g).\]

Because this action distributes over the product of functions, it satisfies the required structural property (a homomorphism) for group actions on algebraic structures.
\end{proof}

\begin{theorem}
The alternating functions are the only unsymmetrical functions of $n$ variables of which a power can be symmetric.
\end{theorem}

\begin{proof}
It is sufficient to prove this theorem for prime powers; for if there exists a function $F(x_1, x_2, \hdots, x_n)$ such that $F^{pq}$ is symmetric and $p$ is a prime, then there is 
also a function $\phi = F^q$ such that $\phi^p$ is symmetric. \\

Therefore let \[ \phi^p = S, \mbox{    a symmetric function.}\] Since the group of $\phi$ which is unsymmetric, cannot contain all the transpositions, let $\sigma = (x_\alpha x_\beta)$ be a 
transposition which converts $\phi$ into $\phi_j$; we have \[ \phi_j^p = \phi^p = S.\] and therefore $\phi_j = \omega\phi$, where $\omega$ is a $p$th root of unity. Hence
\[ \sigma \phi = \phi_j = \omega \phi,\] and operating again with $\sigma$ 
\[ \sigma^2 \phi = \omega\sigma\phi = \omega^2 \phi,\] but $\sigma^2 = I$ and therefore $\omega^2 = 1$, and consequently $p=2$. \\

Since $\phi^s$ is symmetric, $\phi$ is an alternating function. \qed
\end{proof}

\begin{theorem}
For any number, $n$, of independent elements there is no multiple-valied function of which a power is two-valiued when $n>4$; and when $n=3$, or $n=4$, if there is any such power it is
third power.
\end{theorem}

Concentrating on prime numbers, and supposing that $\phi$ is a multiple-valued function whose $p$ power is two-valued we have
\be \phi^p = S_1 + S_2\sqrt{\triangle}. \label{two-valued}\ee
The group of $\phi$ cannot contain all the circular substitutions of the third order, for if it did, then this group would coincide with the alternating group, and $\phi$ would be two-valued. Let
$\sigma = (x_\alpha, x_\beta, x_\gamma)$ be such a substitutions not contained in the group of $\phi$ and suppose $\sigma\phi = \phi_j$. From Eq. (\ref{two-valued}),
since $S_1 + S_2\sqrt{\triangle}$ is unaltered by $\sigma$, we have
\[ \phi^p = \phi_j^p,\] hence $\phi_j = \omega \phi$, where $\omega$ is a $p^{th}$ root of unity. Operating again twice in succession with $\sigma$ we obtain

\begin{eqnarray*}
\sigma \phi &=& \omega \phi,\cr
\sigma^2 \phi &=& \omega\sigma\phi = \omega^2 \phi,\cr
\sigma^3 \phi &=& \omega^2 \sigma \phi = \omega^3 \phi,\cr
\end{eqnarray*}whence, since $\sigma^3 = 1$ (three cycle), we have $\omega^3 = 1$, and therefore $p=3$.\\

Again, when the number of elements is greater than 4, there are circular substitutions of the fifth order, and those cannot all be contained in the group of $\phi$ (or else it will be two-valued and belong to the Alternating Group). Let $\tau$ be one of those not contained in this group and $\tau \phi = \phi_j$. We have, as before, from Eq. (\ref{two-valued}), by applying this substitution
(which does not affect the right-hand side),
\[ \phi^p = \phi_j^p = S_1 + S_2\sqrt{\triangle}.\] Hence proceeding as before, we have $\tau\phi = \omega\phi$ and oeprating again on this and the succeeding equations with $\tau$ we
easily find $\tau^5 \phi = \omega^5 \phi$; whence $\omega^5 = 1$, since $\tau^5 = 1$, and it is proved that $p=5$. Now this results being incompatible with the previous value of $p$
obtained, we infer that when the number of elements is greater than 4, it is impossile to find any multiple-valiued function $\phi$, a prime power of which will be two-valued.\\

There are actually, when $n$ is not greater than 4, multiple-valiued functions, a third power of which is two-valued.